\section{The complete local Type-I theorem}
\label{sec:type-one-closure}

\status{TC1--TC6 complete the source's supplied-data
\texttt{Conclusion2}. The proof uses the actual maps, ranks and
numerical bounds reconstructed above. In particular, no unproved
\texttt{AdmB} or \texttt{AdmT} package is assumed. The existence of
the simultaneous global arithmetic data remains separate.}

\paragraph{TC1: the cloud count is a finite matrix argument.}
Write $G=(p^{n'k}-1)/(p^{n'}-1)$ and
$d(\lambda)=\dim V(\lambda)/G$. HM11 constructs the actual maps
\[
 i:V(H)\longrightarrow V(H^+),\quad
 U_b:V(H)\longrightarrow V(\lambda),\quad
 L_a:V(\lambda)\longrightarrow V(H^+)
\]
with
\[
 L_aU_b=\delta_{ab}p^{\kappa_b}i,\qquad
 0\le\kappa_b\le K,\qquad \#\{b\}=B.                    \tag{TC.1}
\]
The exact midpoint-lattice preimage makes $i$ a primitive
embedding. Extending an integral basis gives a retraction that
preserves the midpoint lattices. Perfect normalized duality
transfers its dual with the factor $p^{A_H-A_{H^+}}$.

Take $n_1,n_2$ good factors for $H,H^+$ at depth $D$, and
absorb the integral column operations so their diagonal scalars
are $z_1,z_2$, of valuation at most $D$. Smith-reduce the cross
pairing matrix $N$ between the image of the first family and
the second dual family. Let $r$ count its invariant factors
dividing $z=z_1z_2$, and put $s=n_1-r$. On the remaining rows
the cross matrix is $zE$ with $E\in wO$.

The combined pairing on these $s$ rows and the second family is
\[
 M_0=\begin{pmatrix}z_1I_s&zE\\ B_0&z_2I_{n_2}\end{pmatrix}.
\]
Both $U=I-EB_0$ and $V=I-B_0E$ are invertible over $O$, and
\[
 M_0
 \begin{pmatrix}
 z_2U^{-1}&-zEV^{-1}\\ -B_0U^{-1}&z_1V^{-1}
 \end{pmatrix}
       =zI.                                             \tag{TC.2}
\]
These are $s+n_2$ factors at depth $2D$ with a cleaned source
and the full midpoint dual. HM5's one-sided upper bound, using
$4D<c_2$, gives $s+n_2\le d(H^+)$. Thus
$r\ge n_1+n_2-d(H^+)$.
The good rows of $N$, followed by (TC.1), give
\[
 Br\ \hbox{factors for }\lambda,\qquad
 r\ge n_1+n_2-d(H^+),\qquad
 \hbox{depth }2D+K.                                    \tag{TC.3}
\]
The normalized pullbacks through $L_a$ have the conductor
$c_2-2Nc$ proved in HM11. Multiplying each by
$p^{K-\kappa_a}$ makes the common scalar exactly $p^Kz$.
This proves the source's cloud count without an additional
intersection-dimension assumption.

\paragraph{TC2: compare the dimensions factor by factor.}
At each selected position let
\[
 \lambda_s=(A_s+c,\beta_{s,1},\ldots,\beta_{s,k-2},-B_s-c),
 \qquad |\beta_{s,i}|\le c,\qquad A_s,B_s\ge b'/3.
\]
Let $H$ have extremes $(A_s,-B_s)$ and zero middle weights;
$H^+$ has extremes $(A_s+2c,-B_s-2c)$. Let $\mathrm{Cl}$ be
the product of the middle Weyl dimensions. The Weyl product CB3
gives
\[
 \mathrm{Cl}\dim H\le\dim\lambda\le E\mathrm{Cl}\dim H,
 \qquad \dim H^+\le E\dim H,\qquad
 E=(1+6c/b')^{N(2k-3)}.                                \tag{TC.4}
\]
Indeed each remaining numerator increases by between zero and
$2c$, or by $4c$ in the extreme-to-extreme factor for $H^+$.
Its corresponding denominator is at least $A_s$, $B_s$ or
$A_s+B_s$.
For $\varepsilon=12Nkc/b'\le1/9$,
$E\le e^\varepsilon$, $(1-\varepsilon)E\le1$ and $E\le2$.
Consequently TC1 converts flat retained fractions
$1-\kappa_1,1-\kappa_2$ into the general fraction
\[
 1-\kappa_1-2\kappa_2-\varepsilon,\qquad
 \hbox{depth }2D+N\lceil\kappa_{\rm one}\rceil,\qquad
 \hbox{conductor }c_2-2Nc.                             \tag{TC.5}
\]
Here $0\le\kappa_1\le1$, $\kappa_2\ge0$.
The sign check is essential:
$(1-\kappa_1)\mathrm{Cl}\,d(H)
\ge(1-\kappa_1)(1-\varepsilon)d(\lambda)$.

\paragraph{TC3: restore the weights outside the selected set.}
Write $\lambda=\lambda_T+\lambda_O$ on disjoint letter sets.
Multiplication in disjoint polynomial variables is an integral
tensor isomorphism for every lattice family. The zero-weight
factors are constants, so dimensions multiply and the normalized
pairing is the tensor product. If all coordinates of $\lambda_O$
are bounded in absolute value by $c$, its first-minus-last weight
sum is at most $h=2n'c$. In an elementary-divisor frame, the
tableau row degrees and the signed determinant twist give
\[
 R_{\lambda_O,e}\subseteq w^{-h}H_{\lambda_O,\rm src},
 \qquad R'_{\lambda_O,e}
       \subseteq w^{-h}H_{\lambda_O,\rm tgt}^{*}.        \tag{TC.6}
\]
Tensor with the exact dual midpoint bases of HM8. The height
increases by $h$, the conductor decreases by $h$, and the depth
is unchanged. Rank multiplies by $\dim V(\lambda_O)$.
Thus the retained \emph{fraction} is unchanged. This proves
the source's \texttt{ResidueClassCount}.

\paragraph{TC4: the edge factors give the actual integral rank.}
There are exactly $G$ incoming rank-one edges at any target
lattice. Choose the local factors on each. Their combined pairing
matrix has diagonal blocks $z_eI$, with $z_e\mid p^D$, and
off-diagonal entries in $w^cO$. After multiplying column $e$
by $p^D/z_e$, the matrix is
\[
                         p^D(I+E),\qquad E\in wO,      \tag{TC.7}
\]
provided $c>2D$. Its remaining inverse is integral. Summing
the local ranks therefore gives at least
$(1-\kappa)\dim V(\lambda)$ factors dividing $p^D$.
The vectors retain their individual other-edge vanishing.
This is the full matrix certificate \texttt{GoodRankK}, hence
\texttt{LPprime}, with the source's stronger conductor hypotheses
$8D<c$ and $4(D+1)\le c$ retained.

\paragraph{TC5: attain the infimum and apply the two shifts.}
Let $\kappa_F(A,B)$ be the least nonnegative flat deficit at
height $t_*(A,B)$, conductor $c_4$ and depth $D_1$.
It lies in $[0,1]$, and its infimum is admissible.
To see this at any edge, put $T=(1-\kappa_F)d(H)$.
If $T>0$, a sufficiently close admissible deficit forces its
integer rank to exceed $\lceil T\rceil-1$, hence to be at least
$T$; if $T\le0$, rank zero suffices.

For $c=c_{\rm cl}$ the actual final deficit is
\[
 \begin{split}
 \kappa_{\rm text}(a_w,b_w)
  ={}&\begin{cases}
      \kappa_F(a_w-c,b_w-c),&
          a_w-c\ge a,\ b_w-c\ge b_l,\\
      1,&\text{otherwise},
      \end{cases}\\
     &+2\kappa_F(a_w+c,b_w+c)+12Nkc/b'.
 \end{split}                                           \tag{TC.8}
\]
All inequalities and shifts are coordinatewise. If this deficit
is at least one, use the zero certificate. Otherwise the first
branch holds. Enlarge its height to $t_*(a_w+c,b_w+c)$ and
apply TC2 to the two flat weights. Apply TC3 next. The resulting
conductor is exactly
\[
 c_4-2Nc-2n'c=c_*.
\]
NP5 and NP7 give respectively
\[
 2D_1+N\lceil\kappa_{\rm one}\rceil\le D_p,\qquad
 t_*(a_w+c,b_w+c)+2n'c\le t_I.                         \tag{TC.9}
\]
TC4 now proves \texttt{MidLP} for every requested integer
height $t\ge t_I$.

\paragraph{TC6: a completely explicit local theorem.}
The following sufficient constants keep all previous estimates:
\[
 \begin{gathered}
 C_0\ge880,\quad C_{E1}=C_{E2}=C_m=17,\quad
 C_\kappa=3,\quad C_t=30,\quad C_c=4,\\
 C_A=50000,\qquad C_s=351,\qquad C_{H,\rm inner}=700000.
 \end{gathered}                                         \tag{TC.10}
\]
NP2 gives the last threshold as $1000\cdot35\cdot5\cdot4$.
The SW mean coefficient is at most
$2(31e^3+5)e^3<45468$ using $e<3$, and NP9 gives
$C_s=3(17+4\cdot17+17+15)=351$.

There is no remaining pointwise good-share assumption.
RU10 makes each segment lattice finite free; the ranks of its
saturated good sublattices form a nonempty finite set. Choose
one of maximal rank. The literal definition of $\kappa_{TI}$
then supplies its good share, including the rank-zero case.
Moreover
\[
 \sum_jD_j+4\kappa_0=D_0,\qquad
 \kappa_P+\sum_j\kappa_{TI,j}=\kappa_x.                 \tag{TC.11}
\]
RU11 therefore supplies exactly HM6's aggregate certificate.
HM6--HM7 and RD11--RD15 give
\[
 \operatorname{mean}\kappa_{\rm text}
     \le \{6(C_A+C_s)+8\}\frac{NN_1}{Y}
     =302114\,\frac{NN_1}{Y}.                          \tag{TC.12}
\]
The final numerical thresholds are
$T_{\rm LP}=400093$, $T_M=2800001$ and
$T_{\rm density}=400093$; the height coefficient only needs
$C\ge124$. Thus we have proved:

\begin{quote}
\textbf{Local Type-I theorem.}
For $C_0\ge880$, $C_H\ge2800001$ and $C=302114$,
the source's literal \texttt{HypI} implies
\texttt{Conclusion2}.
There is a nonnegative exceptional function of the degree
families, with mean at most $CNN_1/Y$, which works for every
allowed coefficient-ring setting and every dominant
\texttt{WeightI}. At every target lattice and every integer
$t\ge t_I$, it gives at least
$(1-\kappa)\dim V(\lambda)$ integral factors dividing $p^{D_p}$,
with the required other-edge vanishing to order $c_*$.
\end{quote}

For uniformity in the coefficient ring, choose one setting and
its $\kappa_{\rm text}$; HM12 transports its \texttt{LPprime}
certificate to every other setting. Equality of two independently
computed exceptional functions is unnecessary. If no setting
exists, use $\kappa=0$ and the conditional assertion is vacuous.
All constants precede the local data. The expanded companion
states every field of \texttt{HypI} and the exact $t_I$ formula.

Together with EA9, this completes both local Type-I and Type-II
constructions. It does not supply the global abundance,
theta/spectral identities, metric comparison, cohomological
windows or simultaneous prime choices still needed for RH
and Dirichlet GRH.
